\honest{Evaporative Cooling of Group Beliefs}{Eliezer Yudkowsky}{2007}

Early students of cults, I say, ``were surprised to discover'' that cults hit by a failed prophecy or a scandal ``often come back stronger than before.'' \nb{The one study the post names, Festinger, Riecken and Schachter's \textsc{When Prophecy Fails} (1956), was not surprised: its authors predicted the effect and joined the group to test it. Whether groups come back stronger is disputed. Melton (1985) held that testing strengthens religious groups; Thomas Kelly argues that failed prophecies are often fatal.} The Jehovah's Witnesses placed Armageddon in 1975, and the Unarians survived a spacefleet that did not arrive. \nb{Both show survival, not a stronger group. By Wikipedia's account, many Witnesses became inactive in 1976--80.}

The usual explanation is cognitive dissonance: people who gave everything for a belief cannot admit they were wrong, so each becomes more fanatical. \nb{The book's account is also social. Its fifth condition is that ``The individual believer must have social support''; an isolated believer, it predicts, will give up.}

I offer another. In evaporative cooling the fastest atoms escape the trap, and the rest grow colder. In \textsc{When Prophecy Fails}, one member walked out right after the saucer failed to land. Who leaves first, I ask: an average member, or a skeptic who had been a brake on the fanatics? \nb{The book agrees. The man who left, Kurt Freund, was ``the sole vocal skeptic in the group,'' and the only two members who gave up the belief were ``lightly committed to begin with.'' The book also reports that the group won no convert and was gone by 9 January. Kelly, drawing on newly unsealed archives, argues that its leader recanted.} Once the skeptics are gone, the debate runs from fanatics to slightly less extreme fanatics. Without its moderates, the group ``may internally increase in average fanaticism.'' \nb{The first step is arithmetic: when the moderates leave, the average of the rest rises with no one changing their mind. None of the three examples tests the second step. Each shows who left.}

When Rand's affair with Branden was revealed, many Objectivists left with Branden for an ``open system'' of Objectivism. \nb{The affair was first made public in 1986. The ``open system'' split came in 1989--90, around David Kelley, not Branden.} On a transhumanist mailing list, a leftist faction insulted the libertarians until they left.

So be slow to eject dissenters. On the other hand, Kuhn thought a science must shut out outsiders before it can do real work, and I can ``see a core of truth here.'' \nb{Kuhn also wrote that those who cling to older views are ``simply read out of the profession,'' which is the ejection the post had just warned against. The post does not say how to tell the two apart.} In a footnote I add that my own theory of moderating online discussions ``may not have served me too well in practice.''
